\documentclass[12pt]{article}

\usepackage[a4paper, margin=1in]{geometry}
\usepackage{amsmath, amssymb, amsthm}
\usepackage{enumitem}
\usepackage{mathtools}

\title{Lecture Scribe -- Probability Models and Random Variables}
\author{}
\date{}

\begin{document}

\maketitle

\section*{1. Multinomial Counting and Partitioning}

\subsection*{Dividing Distinct Objects into Equal Groups}

If 20 distinct objects are divided into 4 groups of 5 objects each, where:
\begin{itemize}
    \item Order within each group does \textbf{not} matter
    \item Order of the groups does \textbf{not} matter
\end{itemize}

The total number of ways is:
\[
\frac{20!}{(5!)^4 4!}
\]

\subsection*{Probability of Specific Grouping}

If each cuisine already has exactly 5 dishes, then the only way each platter contains dishes of a single cuisine is when each cuisine forms its own platter.

Since platters are unlabeled, this arrangement is unique.

\[
P = \frac{1}{\frac{20!}{(5!)^4 4!}} 
= \frac{(5!)^4 4!}{20!}
\]

\section*{2. Uniform Distribution (Continuous Case)}

If arrival time is uniformly distributed over a 30-minute interval:

\[
P(\text{event}) = \frac{\text{favorable time}}{\text{total time}}
\]

\subsection*{Waiting Time Probabilities}

Total interval: 30 minutes

\[
P(\text{wait} < 5) = \frac{10}{30} = \frac{1}{3}
\]

\[
P(\text{wait} > 10) = \frac{10}{30} = \frac{1}{3}
\]

\section*{3. Conditional Probability and De Morgan's Law}

Let:
\[
L: \text{employee arrives late}, \quad
E: \text{employee leaves early}
\]

Given:
\[
P(L)=0.15, \quad P(E)=0.25, \quad P(L \cap E)=0.08
\]

\subsection*{Complement Events}

\[
P(E^c) = 1 - P(E) = 0.75
\]

\subsection*{Union Formula}

\[
P(L \cup E) = P(L) + P(E) - P(L \cap E)
\]

\[
P(L \cup E) = 0.32
\]

\subsection*{De Morgan's Law}

\[
P(L^c \cap E^c) = 1 - P(L \cup E)
\]

\[
P(L^c \cap E^c) = 0.68
\]

\subsection*{Conditional Probability}

\[
P(L^c \mid E^c) = 
\frac{P(L^c \cap E^c)}{P(E^c)} 
= \frac{0.68}{0.75} \approx 0.907
\]

\section*{4. Poisson Distribution}

\subsection*{Definition}

If $X \sim \text{Poisson}(\lambda)$,

\[
P(X = k) = \frac{e^{-\lambda}\lambda^k}{k!}
\]

\subsection*{Example 1: $\lambda = 0.2$}

\[
P(X=0) = e^{-0.2} = 0.8187
\]

Using complement rule:

\[
P(X \ge 2) = 1 - P(0) - P(1)
\]

\[
P(X \ge 2) = 0.0176
\]

\subsection*{Example 2: $\lambda = 3.5$}

\[
P(X \ge 2) = 1 - P(0) - P(1)
\]

\[
P(X \ge 2) = 0.8641
\]

\[
P(X \le 1) = P(0) + P(1) = 0.1359
\]

\subsection*{Example 3: $\lambda = 3$}

\[
P(X \ge 3) = 1 - P(0) - P(1) - P(2)
\]

\[
P(X \ge 3) = 0.5768
\]

Conditional probability:

\[
P(X \ge 3 \mid X \ge 1) 
= \frac{P(X \ge 3)}{P(X \ge 1)} 
\approx 0.607
\]

\section*{5. Gaussian (Normal) Distribution and Q-Function}

Given PDF:

\[
f_X(x) = \frac{1}{\sqrt{8\pi}} 
\exp\left(-\frac{(x+3)^2}{8}\right)
\]

Comparing with Gaussian form:

\[
f_X(x)=\frac{1}{\sqrt{2\pi\sigma^2}} 
\exp\left(-\frac{(x-m)^2}{2\sigma^2}\right)
\]

We identify:

\[
m=-3, \quad \sigma^2=4, \quad \sigma=2
\]

\subsection*{Relations}

\[
F_X(x)=\Phi\left(\frac{x-m}{\sigma}\right)
\]

\[
P(X>x)=Q\left(\frac{x-m}{\sigma}\right)
\]

\subsection*{Key Probability Expressions}

\[
P(X \le 0)=\Phi\left(\frac{3}{2}\right)
=1-Q\left(\frac{3}{2}\right)
\]

\[
P(X>4)=Q\left(\frac{7}{2}\right)
\]

\[
P(|X+3|<2)=1-2Q(1)
\]

\[
P(|X-2|>1)=1-Q(2)+Q(3)
\]

\section*{6. Jury Decision Model (Binomial Model)}

Each juror makes the correct decision independently with probability $\theta$.

Conviction requires at least 8 of 12 jurors.

If $\alpha$ is the probability the defendant is guilty, then probability of correct decision:

\[
\alpha \sum_{i=8}^{12} 
\binom{12}{i}\theta^i(1-\theta)^{12-i}
+
(1-\alpha)
\sum_{i=5}^{12}
\binom{12}{i}\theta^i(1-\theta)^{12-i}
\]

\section*{7. Discrete Distribution and CDF}

Given:

\[
p(i)=c\frac{\lambda^i}{i!}, 
\quad i=0,1,2,\dots
\]

Using:

\[
\sum_{i=0}^{\infty}
\frac{\lambda^i}{i!}=e^\lambda
\]

We obtain:

\[
c=e^{-\lambda}
\]

Thus:

\[
P(X=0)=e^{-\lambda}
\]

\[
P(X>2)=
1-e^{-\lambda}
-\lambda e^{-\lambda}
-\frac{\lambda^2 e^{-\lambda}}{2}
\]

\subsection*{Cumulative Distribution Function (Discrete Case)}

\[
F(a)=\sum_{x\le a}p(x)
\]

For discrete $X$, $F$ is a step function.

Example:

\[
p(1)=\frac14,\quad
p(2)=\frac12,\quad
p(3)=\frac18,\quad
p(4)=\frac18
\]

\[
F(a)=
\begin{cases}
0, & a<1 \\
\frac14, & 1\le a<2 \\
\frac34, & 2\le a<3 \\
\frac78, & 3\le a<4 \\
1, & a\ge4
\end{cases}
\]

\section*{8. Reliability of Systems (Binomial Model)}

If $n$ components function independently with probability $p$, system works if at least half function.

\subsection*{Case 1: 5 vs 3 Components}

5-component system effective probability:

\[
\binom{5}{3}p^3(1-p)^2
+\binom{5}{4}p^4(1-p)
+p^5
\]

3-component system effective probability:

\[
\binom{3}{2}p^2(1-p)
+p^3
\]

5-component system better if:

\[
p>\frac12
\]

\subsection*{General Case}

A $(2k+1)$-component system is better than a $(2k-1)$-component system if and only if:

\[
p>\frac12
\]

Because:

\[
P_{2k+1}-P_{2k-1}
=
\binom{2k-1}{k}
p^k(1-p)^k(2p-1)
\]

Positive if and only if:

\[
2p-1>0
\Rightarrow
p>\frac12
\]

\end{document}